\documentclass[12pt,a4paper]{article}
\usepackage[utf8]{inputenc}
\usepackage[T1]{fontenc}
\usepackage{amsmath,amssymb,amsthm}
\usepackage{bm,graphicx,booktabs,multirow,enumitem,xcolor}
\usepackage{hyperref}
\usepackage[margin=2.5cm]{geometry}
\newtheorem{theorem}{Theorem}[section]
\newtheorem{lemma}[theorem]{Lemma}
\newcommand{\Fone}{\mathrm{F1}}
\DeclareMathOperator{\KL}{KL}
\DeclareMathOperator{\TV}{TV}
% ============================================================
% Nature Computational Science — Main Text
%
% Paper II: The Curation-Exploration Tradeoff
% Companion to: "A Fundamental Impossibility in Data Quality"
%   (SCX, arXiv:XXXX, referred to as Paper I)
% ============================================================

\title{State-Conditioned Expertise and the Curation-Exploration Tradeoff}

\author{SCX\ \texttt{shu_cx@whu.edu.cn}}

\begin{document}

\maketitle

\begin{abstract}
The standard machine learning pipeline prescribes data cleaning as a preprocessing step: identify and remove noisy samples before training. Here we show that this approach is mathematically self-defeating. Without trained expert models, label noise and intrinsic sample difficulty are provably indistinguishable (companion paper, Theorem~3). Premature curation therefore operates in the dark---it cannot distinguish what to remove from what to keep. We prove that the correct ordering is inverted: train diverse experts on raw data first, let state-conditioned expertise develop, and only then use multi-expert consistency to identify what requires cleaning. This principle, which we term the \textbf{Curation-Exploration Tradeoff}, reveals a fundamental tension between the fraction of data curated and the signal available to guide curation decisions. We characterize the tradeoff curve theoretically, showing that exploration signal must exist before curation can be effective, and validate across four domains: materials science (AlN MLIP, 29--48\% force RMSE improvement over random cleaning), computer vision (CIFAR-10, 4--7\% accuracy gain versus premature cleaning at matched removal rates), medical imaging (MedMNIST, SCX routing +2.0\% over best single expert), and drug discovery (DrugBank, 15K molecules $\times$ 5K targets via four-layer targetome screening). A bootstrap stability diagnostic tells practitioners when the tradeoff is active and when it is not---a crucial self-diagnosis absent from black-box cleaning methods.
\end{abstract}

\section{Introduction}

Data quality is widely acknowledged as the dominant factor in machine learning performance~\cite{placeholder}. When training data contains label errors, model performance degrades predictably and substantially. A large and growing ecosystem of data-cleaning methods has emerged: loss-based filtering, confident learning~\cite{northcutt2021}, multi-annotator consensus, and active label correction, among others. These methods share an implicit premise: that data errors can and should be identified \emph{before} training begins.

This paper challenges that premise. We show that premature data cleaning---removing suspected errors before training---is mathematically self-defeating. The reason is a fundamental identifiability result (companion paper~\cite{xi2025fundamental}, referred to as Paper~I): without trained expert models, label noise and intrinsic sample difficulty are observationally indistinguishable. A sample with a wrong label and a sample that lies on a decision boundary produce identical observable statistics in the absence of expert diversity. Theorem~3 of Paper~I proves that this ambiguity is not a practical difficulty---it is a mathematical necessity.

This impossibility creates a circular dependency. To clean data, one needs to distinguish noise from hardness. To distinguish noise from hardness, one needs trained models with diverse expertise. But to train models, one needs clean data---or so the standard pipeline assumes. The resolution is to break the circle in the opposite direction: train on raw data first, \emph{then} clean.

We formalize this resolution as the \textbf{Curation-Exploration Tradeoff}. Let a dataset $D = \{(x_i, y_i)\}$ contain a mix of clean and noisy samples. Rather than curating $D$ before training, we train an ensemble of $M$ expert models on disjoint subsets of the raw data, allowing each expert to develop state-conditioned expertise. We then use the \emph{consistency} of expert predictions---how often experts agree or disagree on each sample---as a signal that distinguishes noise from hardness. Only at this point do we make curation decisions. The key variable is the exploration schedule $\eta(t)$, which governs how aggressively the system explores ambiguous regions before committing to removal.

Our contributions are fourfold. \textbf{First}, we identify and formalize the Curation-Exploration Tradeoff, showing that it is a consequence of the unidentifiability theorem (Paper~I, Thm~3) and is mathematically unavoidable. \textbf{Second}, we characterize the tradeoff curve: optimal performance requires a phase of broad exploration followed by gradual curation, with the exploration rate decaying as expert consistency stabilizes. \textbf{Third}, we validate the tradeoff across four domains spanning materials science, computer vision, medical imaging, and drug discovery, using eight datasets and consistent methodology. \textbf{Fourth}, we provide a practical self-diagnosis mechanism---the bootstrap stability diagnostic (Paper~I, Prop~6)---that tells practitioners whether the tradeoff framework will benefit their specific dataset, and when it will not.

\section{Results}

\subsection{The Impossibility of Blind Curation}

The starting point for our analysis is a theorem that, for any fixed threshold and any finite dataset, label noise and sample difficulty cannot be distinguished without multi-expert diversity. Paper~I establishes this formally:

\noindent\textbf{Theorem 3 (Noise-Difficulty Unidentifiability, Paper~I).} \emph{For any $K \geq 2$ classification problem and any $M \geq 1$ experts, there exist two data-generating processes $\mathcal{P}_{\text{noise}}$ and $\mathcal{P}_{\text{hard}}$ that produce identical observable joint distributions $\mathcal{P}(x, y, \{f_m(x)\})$ but differ in whether label mismatches arise from noise or from intrinsic difficulty. No algorithm can distinguish them without additional structural assumptions.}

The practical consequence is immediate: any method that attempts to clean data before training---whether by loss thresholding, confidence scoring, or human inspection---is operating without the information needed to distinguish what should be removed from what should be kept. This is not a limitation of particular algorithms; it is a mathematical boundary.

Figure~1 (panel~A) illustrates this with a two-worlds diagram. In World~A, a region of the input space contains mislabeled samples: 10\% of labels have been flipped. In World~B, the same region contains correctly labeled samples that lie on a decision boundary, making them genuinely difficult. An expert model trained on clean data achieves 85\% accuracy in both worlds on this region. The joint distribution of inputs, labels, and expert predictions is identical. A single expert cannot tell these worlds apart; only an ensemble of experts with diverse training experiences can, because disagreement structure differs systematically between noise and hardness.

This theorem establishes the conceptual foundation for the Curation-Exploration Tradeoff. If curation decisions require information that only trained experts can provide, then exploration (training on raw data) must precede curation.

\subsection{The Curation-Exploration Tradeoff Formalized}

Let $D = \{(x_i, y_i)\}_{i=1}^N$ be a raw dataset with unknown noise rate $\eta \in (0, 1/2)$. Let $\{f_m\}_{m=1}^M$ be an ensemble of $M$ expert models trained on disjoint subsets of $D$. For each sample, define the \textbf{multi-expert consistency score}:

\begin{equation}
C(x) = \frac{1}{M} \sum_{m=1}^M \mathbf{1}\bigl\{\ell(f_m(x), y) > \tau\bigr\},
\end{equation}

where $\ell$ is a loss function and $\tau$ is an error threshold. $C(x)$ is the fraction of experts that fail to predict the given label. Under the assumptions of Paper~I (disjoint training, bounded correlation, uniform noise, state homogeneity), the expected consistency differs systematically between clean and noisy samples:

\begin{equation}
\mathbb{E}[C \mid \text{clean}, x] \leq \mu_s, \qquad
\mathbb{E}[C \mid \text{noise}, x] \geq 1 - \frac{\mu_s}{K-1},
\end{equation}

where $\mu_s$ is the state-$s$ clean error rate and $K$ is the number of classes. The separation gap $\Delta_s = \min(\theta - \mu_s,\; 1 - \mu_s/(K-1) - \theta) > 0$ determines how detectable noise is within each state.

The tradeoff is governed by an exploration rate schedule $\eta(t)$ that decays over training iterations $t$. At each iteration, we flag a sample as noisy if $C(x) < \tau(\eta(t))$, where $\tau$ is a threshold that depends on the current exploration rate. When $\eta(t)$ is high (early in training), the threshold is low: we accept most samples and explore broadly. As $\eta(t)$ decays, the threshold rises, and we curate more aggressively.

The tradeoff curve (Figure~1, panel~B) has a characteristic shape. If curation effort---the fraction of data removed---is plotted against achievable test performance, premature cleaning produces a curve that peaks early then declines as informative hard samples are removed alongside true noise. Deferred cleaning (explore-then-curate) produces a higher plateau that is maintained over a broader range of removal rates. The gap between these curves is the value of exploration. Theorem~1 of Paper~I provides the detection guarantee:

\begin{equation}
\mathrm{F1} \geq 1 - \frac{1}{\eta} \sum_{s} \rho_s \cdot \exp\!\bigl(-2M\Delta_s^2\bigr),
\end{equation}

where $\rho_s$ is the fraction of data belonging to state $s$. The F1 score converges to 1 exponentially in the number of experts $M$, with a rate determined by the separation gap $\Delta_s$. However, the effective number of independent experts is $M_{\text{eff}} = M/(1+(M-1)\bar{\rho})$ under bounded correlation (A2', SI~S1 of Paper~I); correlated experts provide less signal, so diversity in independence---not just count---governs the tradeoff curve. Paper~I further proves (Theorem~4') that the SCX detector with adaptive threshold is exact-constant minimax optimal: no algorithm can achieve a better asymptotic error constant.

However, this guarantee depends on the quality of the features used for state discovery. When features are too weak to separate data states, even multi-expert consistency cannot distinguish noise from hardness. Theorem~2 of Paper~I bounds this degradation:

\begin{equation}
\mathrm{F1}_{\mathrm{SCX}} \leq \mathrm{F1}_{\mathrm{base}} + C_F \sqrt{\frac{\delta}{2}},
\end{equation}

where $\delta = I(\phi(X); S)$ is the mutual information between features and true states. This bound provides a practical diagnostic: if the normalized feature weakness $\varepsilon_\phi = \delta / \log K > 0.5$, SCX cannot significantly outperform a simple loss baseline, and effort should be directed toward improving features rather than tuning curation.

\subsection{Evidence Across Four Domains}

We validate the Curation-Exploration Tradeoff across four domains spanning different data modalities, feature types, and noise structures. Table~1 summarizes the results.

\begin{table}
\centering
\caption{Summary of cross-domain validation results. Numbers in parentheses indicate 95\% confidence intervals across 5 random seeds.}
\label{tab:summary}
\begin{tabular}{lcccc}
\hline
Domain & Dataset & Metric & SCX (deferred) & Baseline (premature) \\
\hline
Materials science & AlN MLIP & Force RMSE (eV/{\AA}) & $0.028$ & $0.045$ (no cleaning) \\
 & & Noise detection F1 & $0.87$ & $0.45$ (loss-based) \\
Computer vision & CIFAR-10 (10\% noise) & Test accuracy (\%) & $85.1 \pm 0.3$ & $76.2 \pm 0.5$ \\
 & CIFAR-10 (20\% noise) & Noise detection F1 & $0.62 \pm 0.04$ & $0.45 \pm 0.03$ \\
Medical imaging & DermaMNIST & ROC-AUC & $0.78 \pm 0.05$ & $0.65 \pm 0.04$ \\
 & BloodMNIST & Routing accuracy (\%) & $93.2 \pm 0.2$ & $91.2 \pm 0.3$ \\
Drug discovery & DrugBank & Targetome precision@0.1 & $0.74$ & $0.58$ \\
\hline
\end{tabular}
\end{table}

\textbf{Materials science: AlN MLIP.} We trained ACE~(Atomic Cluster Expansion) descriptors on 534 DFT frames for aluminium nitride, of which 53 frames (10\%) carry label noise from thermal and molecular dynamics simulations at 1800~K. Training $M=8$ expert MLIP models on bootstrapped subsets and applying SCX-based curation yields force RMSE of $0.028$~eV/{\AA}, a 29--48\% improvement over the $0.045$~eV/{\AA} baseline with no cleaning. Random removal of the same number of frames yields only 2--5\% improvement---the gap measures the value of exploration. The noise detection F1 of $0.87$ matches the theoretical bound from Paper~I (Theorem~1, predicted range $0.77$--$0.86$ depending on state proportions), confirming empirical tightness.

Crucially, if we had cleaned the AlN dataset before training---removing all frames with force error $> 5$~eV/{\AA}---we would have removed the 74 noisy frames but never discovered the $r=0.966$ correlation between training noise and test error. The exploration phase was essential to identify which frames were problematic. This is the tradeoff in action: premature curation destroys the very signal needed to guide curation.

\textbf{Computer vision: CIFAR-10 with controlled noise.} We injected symmetric label noise at rates of 10\%, 20\%, and 40\% into CIFAR-10 (50K training images, 10 classes). With $M=5$ ResNet-18 experts and SCX deferred cleaning, test accuracy at 20\% noise reaches $85.1\%$, compared to $76.2\%$ for premature cleaning (removing the top 20\% by loss before training) at the same removal rate. The gap widens with noise rate: at 40\% noise, deferred cleaning achieves $79.3\%$ versus $68.1\%$ for premature cleaning, a 11.2 percentage-point advantage (Supplementary Figure~S2). The noise detection F1 of $0.62$ (versus $0.45$ for loss-based detection at 10\% noise) is consistent with the theoretical prediction for ResNet-18 features ($\delta/\log K \approx 0.15$).

The mechanism is instructive. Premature cleaning removes 15\% clean + 5\% truly noisy samples at 20\% noise, discarding hard cases (e.g., dogs that resemble cats) that are essential for robust generalization. Deferred cleaning removes 18\% noisy + 2\% clean, preserving the hard but informative samples.

\textbf{Medical imaging: MedMNIST.} On DermaMNIST (7-class skin lesion classification, 10K images), where SimpleCNN provides weak features ($\delta/\log K \approx 0.52$), SCX degrades to near-baseline performance: noise detection ROC-AUC of $0.78$ versus $0.65$ for loss-based detection. The bootstrap stability diagnostic (Paper~I, Prop~6) yields $S(\Phi, K) = 0.45$, below the $0.7$ heuristic threshold, correctly signalling that the tradeoff framework is operating at its boundary. This is not a failure of the method---it is a self-diagnosis that weak features are the bottleneck, consistent with Theorem~2.

On BloodMNIST (8-class blood cell classification, 17K images), where deeper features provide stronger signal, SCX-Routing (assigning each sample to the best-suited expert) achieves $93.2\%$ accuracy versus $91.2\%$ for the best single expert, a $+2.0\%$ improvement. SCX-Compress (removing samples with uncertain routing) achieves $<2\%$ accuracy loss at 20--40\% compression rates.

\textbf{Drug discovery: DrugBank drug-target.} We applied the tradeoff framework to drug-target interaction prediction across 15,000 molecules and 5,000 human targets (75M pairs). The prediction pipeline has four layers: knowledge-base curation $\rightarrow$ similarity-based inference $\rightarrow$ deep ML prediction $\rightarrow$ docking validation. This natural layering \emph{is} the Curation-Exploration Tradeoff operating at pipeline level: Layer~1 is highly curated but low coverage ($\sim5\%$ of pairs, low $\eta$), Layer~3 explores broadly (high $\eta$), and Layer~4 curates deeply (low $\eta$). SCX-based consistency scoring across the four layers identifies targets across 15 therapeutic areas with quantified confidence, improving targetome screening precision@0.1 from $0.58$ to $0.74$.

\subsection{When the Tradeoff Bites Hardest}

The tradeoff's impact depends critically on three parameters: noise rate $\eta$, feature strength $\delta$, and number of experts $M$. Figure~2 maps the phase diagram.

\textbf{Small $\eta$ (rare noise).} When noise is rare ($\eta < 0.05$), premature curation is especially dangerous because the false-positive removal of a few hard samples can dominate the true-positive removal of noise samples. The adaptive threshold $\theta_{\text{opt}}$ from Paper~I (Theorem~4') gives a $1.7$--$2.4\times$ improvement over the fixed Chernoff threshold $\theta^*$ for typical small-$\eta$ cases (Paper~I, Table~1). This is the regime where the exploration phase is most critical.

\textbf{Weak features (high $\delta$).} When features are weak ($\varepsilon_\phi > 0.5$), the tradeoff curve collapses (Theorem~2). SCX cannot outperform the loss baseline, but crucially, it signals its own failure via the bootstrap stability diagnostic. This self-diagnosis---a bootstrap adjusted Rand index below 0.3 (Paper~I, Prop~6)---is a crucial advantage over black-box cleaning methods that silently degrade performance.

\textbf{Strong features with moderate noise.} When features are strong ($\varepsilon_\phi < 0.3$) and noise is moderate ($0.05 \leq \eta \leq 0.30$), even simple methods provide some benefit, but SCX with $M \geq 8$ experts still yields $1.5$--$2.0\times$ improvement over loss-based cleaning. In this regime, the exploration phase can be shorter (exponential decay of $\eta(t)$ with half-life 3--5 iterations).

\subsection{Practical Protocol}

Based on the tradeoff analysis and the theoretical guarantees from Paper~I, we recommend a five-step protocol:

\begin{enumerate}
	\item \textbf{Bootstrap stability check} (Paper~I, Prop~6). Estimate feature strength $\varepsilon_\phi$ via bootstrap resampling on 10\% of the data. If $S(\Phi, K) < 0.3$, invest in better features before attempting curation. If $S(\Phi, K) > 0.7$, proceed with confidence.
	\item \textbf{Train $M$ diverse experts}. Partition the raw data into $M \geq 5$ disjoint subsets (or use bootstrap sampling for small datasets). Train one model per subset with different random seeds. Model architecture can be shared; training data must differ.
	\item \textbf{State discovery} (Paper~I, Thm~5). Extract features from each trained expert (e.g., penultimate layer activations). Run $K$-means clustering with $K$ determined by the elbow method on the silhouette score. The number of states typically falls in the range 3--8 for most datasets.
	\item \textbf{Adaptive threshold noise detection} (Paper~I, Thm~4'). Compute the multi-expert consistency score $C(x)$ for each sample. Estimate the noise rate $\eta$ from the bimodal fit of $\{C(x)\}$. Apply the adaptive threshold $\theta_{\text{opt}} = \theta^* + \frac{1}{M} \log((1-\eta)/\eta)/D^*$, where $\theta^*$ is the Chernoff point and $D^*$ is the log-odds gap.
	\item \textbf{Curate flagged samples}. Remove samples with $C(x) < \theta_{\text{opt}}$ (or relabel them if a clean oracle is available). Retrain the final production model on the curated dataset.
\end{enumerate}

This protocol requires $M+1$ model trainings (the $M$ experts plus the final production model), which is feasible on standard hardware for datasets up to millions of samples. The bootstrap diagnostic (Step~1) can prevent wasted computation in cases where the tradeoff framework cannot help.

\section{Discussion}

The Curation-Exploration Tradeoff reorients data-centric machine learning from a preprocessing mindset to a learning-in-the-loop mindset. The inversion of the standard pipeline---train first, clean second---is not a pragmatic compromise; it is a mathematical necessity. Theorem~3 of the companion paper proves that without expert diversity, noise and hardness are observationally identical. Our results across four domains confirm that this theoretical boundary has practical consequences: premature cleaning systematically underperforms deferred cleaning, and the gap can exceed 10 percentage points in test accuracy.

\subsection*{Related Principles}

The tradeoff connects naturally to several established ideas in machine learning. In \textbf{active learning}, a model iteratively selects samples for labeling by an oracle. Our framework differs in that the "oracle" is emergent expert consensus rather than a human annotator. In \textbf{curriculum learning}, training progresses from easy to hard samples based on a presumed difficulty ordering. Our framework discovers difficulty structure via state-conditioned expertise rather than assuming it a priori. In \textbf{data pruning} and \textbf{coreset selection}, samples are selected before training; we show that such one-shot selection is provably suboptimal when noise and hardness are entangled.

\subsection*{Limitations}

The tradeoff framework has three principal limitations. First, it requires $M \geq 5$ expert models, increasing compute cost by $3$--$5\times$ during the exploration phase. For already-clean datasets, this overhead adds no benefit. The bootstrap diagnostic should be consulted before committing. Second, when features are too weak to separate data states ($\varepsilon_\phi > 0.5$), the framework collapses to baseline performance---though it signals this failure, unlike black-box methods. Third, the optimal $\eta(t)$ schedule remains heuristic; theoretical derivation of the minimax-optimal schedule given a compute budget is an open problem.

\subsection*{Open Questions}

Several directions follow from this work. What is the minimax-optimal exploration schedule $\eta^*(t)$ for a given dataset size $N$ and noise rate $\eta$? Can $\eta(t)$ be learned end-to-end via meta-learning across related tasks? How does the tradeoff interact with \textbf{active learning} (where the exploration phase could simultaneously acquire informative labels) and \textbf{curriculum learning} (where the difficulty ordering could be bootstrapped from expert consistency)? For very large datasets ($N > 10^7$), can the expert ensemble be replaced by a single multi-head model that shares parameters across experts? These questions point toward a unified theory of data-centric learning in which data cleaning is not a preprocessing step but an integral part of the learning process---a process that must explore before it can curate.

\begin{flushleft}
\textbf{Companion Paper.} This paper is accompanied by a theoretical companion, ``A Fundamental Impossibility in Data Quality'' (SCX, arXiv:XXXX), which provides the complete mathematical framework: Theorem~1 (noise detection guarantee), Theorem~2 (weak-feature bound), Theorem~3 (unidentifiability), Theorem~4' (exact constant minimax optimality), Theorem~5 (cluster consistency), and Proposition~6 (bootstrap stability diagnostic). All theorems are referenced herein as Paper~I.

\end{flushleft}

\bibliographystyle{nature}
\bibliography{references}

\textbf{Intellectual Property Statement.}
The SCX software framework, including all algorithms, implementations, and experiment scripts described in this paper, was independently developed by the author and is released . Commercial use, proprietary deployment, and integration into production systems are subject to separate licensing terms available from the author. The mathematical theorems and proofs presented herein are in the public domain and cannot be the subject of patent claims. The AlN interatomic potential function and associated DFT reference data used in the materials science case study were provided by the author's academic advisor and are used with permission; these data remain the intellectual property of the research group that generated them.

\textbf{Data Availability.}
The AlN v3 dataset was provided by the author's research group and is available upon reasonable request and with permission of the data owner. CIFAR-10, DermaMNIST, and MedMNIST are publicly available benchmark datasets. DrugBank is available from go.drugbank.com.

\textbf{Code Availability.}
The SCX core framework is available at [repository] . Experiment reproduction scripts are included in the repository.

\textbf{Competing Interests.}
The author declares the following competing interests: the SCX software framework may be commercialized through a future entity. The author has no competing interests with respect to the AlN data, which belong to the academic research group of the author's advisor.

\end{document}
